\documentclass[12pt]{article}
\usepackage{amsmath,amssymb}

\begin{document}
\section*{{2025 AIME I Problems/Problem 14}}
Source: AoPS Wiki

\subsection*{Solution}
Assume $AX=a$ , $BX=b$ , $CX=c$ , $DX=d$ , $EX=e$ , by Ptolemy inequality we have $a+2d\geq \sqrt{3}XE; a+2c\geq \sqrt{3}BX$ , and as we are minimizing the value of the LHS, we want the inequality reached; this happens when both quadrilaterals $CXAB$ and $AXDE$ are concyclic. We notice that we have right angles $\angle{BXA}=\angle{BCA}$ and $\angle{EDA}=\angle{EXA}$ as they both subtend a diameter of circles by noticing that we have formed Pythagorean triples (along with being told that $\angle{B}=\angle{E}=60^{\circ}$ . And so, since we have supplementary angles, $B$ , $X$ , and $E$ are collinear. Thus, we add up the two equalities to get $2a+2c+2d=\sqrt{3}(XE+BX)$ , or $a+c+d=\frac{\sqrt{3}}{2}(BE)$ . Thus we want to find $a+b+c+d+e=(1+\frac{\sqrt{3}}{2})BE$ . By Law of Cosines, we know that $\cos(\angle{DAC})=\frac{1}{7}$ , and since $\angle{CAB}=\angle{DAE}=30^{\circ}$ , we know $\cos(\angle{EAB})=cos(\angle{CAB}+\angle{DAE}+\arccos(\frac{1}{7}))=cos(60^{\circ}+\arccos(\frac{1}{7}))$ , which we find with the cosine angle addition formula and the fact that $\sin(\arccos(\frac{1}{7}))=\sqrt{1-(\frac{1}{7})^2}$ . Then, once again with Law of Cosines, we find that $BE=38$ , and thus the minimum of $a+b+c+d+e$ , or $(1+\frac{\sqrt{3}}{2})BE$ , is $38+19\sqrt{3}\implies\boxed{060}$ . ~Bluesoul ~hashbrown2009 Rewritten and edited by Juwushu .

\end{document}
